% ====================
% 3. Agent and task setup
% ====================
\section{Agent and Task Setup}
\label{sec:experimental-setup}

\subsection{Implementation and Model Configuration}

The agent is
\texttt{qwen/qwen3.8-flash} through OpenRouter, with
temperature 0, seed 4213, a 1,024-token output limit,
and reasoning disabled. The agent prompt template and
sampling settings are fixed across conditions, with
retrieval limited to ten results.

Qwen was selected after an operational pilot: it
reported zero reasoning usage on 40 pilot calls,
whereas DeepSeek returned empty text on three of
eight calls when reasoning was disabled.

Both memory systems use a common interface and fresh
namespaces for each episode and condition. Mem0 uses
its default extraction and search with a local embedder
and vector retrieval without BM25; Supermemory uses its
hosted API and waits for ingestion before retrieval
(details in Appendix~\ref{app:reproducibility}).


\subsection{Experimental Conditions and Baselines}
\label{sec:conditions}

We evaluate eight conditions: no memory (current probe only),
full transcript (all earlier sessions), Mem0 and Supermemory
with native retrieval, and write-time (pre-extracted, tagged
memories) and read-time (retrieved memories tagged afterward)
scope-tagged variants of each memory system.

Each condition runs all 40 episodes (320 runs). The scope-layer
conditions change tags, resolution rules, and memory ordering
together, so their effects cannot be isolated.

